\section{Key Lifecycle}


\subsection{Generation}

Key generation produces three artifacts: a secret key, a public evaluation key (bootstrap key), and optionally a key-switching key:

\begin{lstlisting}[language=Go]
sk, bsk, err := fhe.GenerateKeys(params)
if err != nil {
    return fmt.Errorf("keygen: %w", err)
}
enc := fhe.NewEncryptor(params, sk)
eval := fhe.NewEvaluator(params, bsk)
dec := fhe.NewDecryptor(params, sk)
\end{lstlisting}

The evaluator does not require the secret key. This separation is the fundamental safety property: the entity performing computation never sees plaintext.

\subsection{Serialization}

All key types implement \texttt{encoding.BinaryMarshaler} and \texttt{encoding.BinaryUnmarshaler}:

\begin{lstlisting}[language=Go]
data, err := sk.MarshalBinary()
// store securely via KMS

var sk2 fhe.SecretKey
err = sk2.UnmarshalBinary(data)
\end{lstlisting}

\subsection{Threshold Key Distribution}

For threshold mode ($t$-of-$n$), key generation uses LSSS (Linear Secret Sharing Scheme):

\begin{lstlisting}[language=Go]
shares, bsk, err := fhe.GenerateThresholdKeys(
    params, t, n,
)
// Distribute shares[i] to party i
// bsk (bootstrap key) is public
\end{lstlisting}

Resharing to a new $(t', n')$ configuration does not require reconstructing the original secret key:

\begin{lstlisting}[language=Go]
newShares, err := fhe.Reshare(
    shares, newT, newN,
)
\end{lstlisting}

